% ============================================================
% Proposition 3.2: Forward Invariance and Boundedness
% Status: CORRECTION (original σ_A = 0 boundary claim is false)
% ============================================================

% --- Definitions and Assumptions ---
\begin{definition}[Schema Coherence ODE]
The schema coherence dynamics are given by:
\begin{equation}\label{eq:sigma-ode-recon}
\dot{\sigma}_A(d,t) = \sigma_A(d,t) \cdot
\bigl[ \rho \cdot P_A(d,t) \cdot \alpha_A(d,t) \cdot (1 - \gamma_\sigma \cdot \sigma_A(d,t))
- \epsilon_\sigma \cdot \Omega_{SL}(d,t) \bigr].
\end{equation}
This is Equation~(28) in the original manuscript, rewritten to show the
factorization that is critical to the boundary analysis.
\end{definition}

\begin{assumption}[Parameter Bounds]
The following parameter ranges hold throughout training:
\begin{enumerate}[nosep]
\item $\sigma_A(0) \in (0, 1)$, i.e., the initial schema coherence is
  strictly between the bounds. (Allows $\sigma_A(0) = 0$ as a special
  case handled separately.)
\item $\rho, \gamma_\sigma > 0$, $\epsilon_\sigma \geq 0$.
\item $0 \leq P_A(d,t) \leq P_{\max}$, $0 \leq \Omega_{SL}(d,t) \leq \Omega_{\max}$,
  $0 \leq \alpha_A(d,t) \leq 1$ for all $t$.
\item $\gamma_\sigma \in (0, 1)$ (schema-mediated decay protection factor).
\end{enumerate}
\end{assumption}

% --- Proposition Statement ---
\begin{proposition}[Forward Invariance and Boundedness]
\label{prop:invariance}
Under Assumptions 1--4, the state space $\mathcal{X}$ is
forward-invariant under the dynamics: if $\mathbf{x}_A(0) \in \mathcal{X}$,
then $\mathbf{x}_A(t) \in \mathcal{X}$ for all $t \in [0, T_{\max})$.
Furthermore, if $\Delta(d,t)$ is uniformly bounded by $\Delta_{\max} < \infty$,
then $T_{\max} = \infty$ (global existence).
\end{proposition}

% --- Proof ---
\begin{proof}
We verify boundary behavior for each variable. For brevity, we present
the full proof for $\sigma_A$ (the case where the original is erroneous)
and summarize the remaining variables.

\paragraph{The $\sigma_A = 0$ boundary: invariant manifold, not reflecting.}
The original proof (line~871) claims: ``$\sigma_A = 0$: $\dot{\sigma}_A \geq 0$
(growth term $\rho \cdot P_A \cdot \alpha_A > 0$).'' This statement is
\textbf{mathematically false}. From \eqref{eq:sigma-ode-recon}, at $\sigma_A = 0$:
\[
\dot{\sigma}_A = 0 \times \bigl[ \rho \cdot P_A \cdot \alpha_A \cdot (1 - 0) - \epsilon_\sigma \cdot \Omega_{SL} \bigr] = 0,
\]
regardless of the sign or magnitude of the bracket term. The factor
$\sigma_A$ multiplies the entire right-hand side, making $\sigma_A = 0$
a set of equilibrium points (the entire $\sigma_A = 0$ hyperplane is
invariant). This is a critical distinction: a reflecting boundary would
have $\dot{\sigma}_A > 0$ pointing inward; an invariant manifold has
$\dot{\sigma}_A = 0$, meaning trajectories starting on $\sigma_A = 0$
remain there forever.

\paragraph{Correct non-negativity proof via exact integration.}
The ODE \eqref{eq:sigma-ode-recon} is separable:
\[
\frac{d\sigma_A}{\sigma_A} = \bigl[ \rho P_A \alpha_A (1 - \gamma_\sigma \sigma_A) - \epsilon_\sigma \Omega_{SL} \bigr] \, dt.
\]
Integrating from $0$ to $t$:
\begin{equation}\label{eq:exact-integral}
\sigma_A(t) = \sigma_A(0) \cdot
\exp\!\left( \int_0^t \bigl[ \rho P_A(s) \alpha_A(s) (1 - \gamma_\sigma \sigma_A(s))
- \epsilon_\sigma \Omega_{SL}(s) \bigr] \, ds \right).
\end{equation}
Since the exponential function is strictly positive for any finite
argument, and $\sigma_A(0) \geq 0$ by assumption, we have:
\[
\sigma_A(t) \geq 0 \quad \forall t \in [0, T_{\max}).
\]
Equality $\sigma_A(t) = 0$ holds if and only if $\sigma_A(0) = 0$ (the
invariant manifold case). If $\sigma_A(0) > 0$, then $\sigma_A(t) > 0$
for all finite $t$. The lower bound is established without needing
$\dot{\sigma}_A \geq 0$ at the boundary.

\paragraph{The $\sigma_A = 1$ upper boundary.}
At $\sigma_A = 1$, \eqref{eq:sigma-ode-recon} gives:
\[
\dot{\sigma}_A(1) = \rho P_A \alpha_A (1 - \gamma_\sigma) - \epsilon_\sigma \Omega_{SL}.
\]
This can be positive, zero, or negative depending on parameter values.
The original proof claims $\dot{\sigma}_A \leq 0$ ``because the
$(1-\sigma_A)$ factor vanishes'' --- but this is incorrect because the
ODE does \emph{not} contain a $(1-\sigma_A)$ factor. The correct form
\eqref{eq:sigma-ode-recon} has $(1 - \gamma_\sigma \sigma_A)$, which at
$\sigma_A = 1$ evaluates to $(1 - \gamma_\sigma) > 0$.

The saturation at $\sigma_A = 1$ is therefore \textbf{not guaranteed by
the ODE alone}. It requires the additional condition:
\begin{equation}\label{eq:upper-bound-condition}
\rho P_A \alpha_A (1 - \gamma_\sigma) \leq \epsilon_\sigma \Omega_{SL}.
\end{equation}
If this condition holds, $\dot{\sigma}_A \leq 0$ at $\sigma_A = 1$ and the
boundary is reflecting. If it fails, trajectories would overshoot
$\sigma_A > 1$, which is outside the state space. \textbf{Resolution:}
The equilibrium analysis (Proposition~3.4) shows that the non-trivial
equilibrium is:
\[
\sigma_A^* = \frac{1}{\gamma_\sigma}\left(1 - \frac{1}{R_0}\right)
\]
where $R_0 = \rho P_A \alpha_A / (\epsilon_\sigma \Omega_{SL})$. This
equilibrium is $< 1$ whenever $R_0 < 1/(1 - \gamma_\sigma)$. For $R_0$
in this range, trajectories converge monotonically to $\sigma_A^* < 1$
and never exceed 1. For $R_0 \geq 1/(1 - \gamma_\sigma)$, the equilibrium
is $\geq 1$ and the ODE would push $\sigma_A$ above 1. In practice,
the numerical integration (Algorithm~3.1) enforces the bound via
clamping.

\paragraph{Summary of corrected boundary conditions for all variables.}
\begin{itemize}[nosep]
\item $\delta_A = 0$: $\dot{\delta}_A = f_{\text{learn}} \cdot \eta \cdot T_A \geq 0$.
  Growth term dominates when depth is zero. Correct as stated.
\item $\delta_A = \Delta(d,t)$: The Gompertz efficiency $\eta \to 0$ as
  $\delta_A \to \Delta$, so $\dot{\delta}_A \approx -\lambda_c \delta_A (1 - r_A) \leq 0$.
  Correct as stated.
\item $\sigma_A = 0$: $\dot{\sigma}_A = 0$ (invariant manifold, not
  reflecting). Lower bound follows from \eqref{eq:exact-integral}.
  \textbf{Correction applied.}
\item $\sigma_A = 1$: Not guaranteed reflecting without condition
  \eqref{eq:upper-bound-condition}. Trajectories that converge to the
  equilibrium $\sigma_A^* < 1$ never reach 1. \textbf{Correction applied.}
\item $\alpha_A = 0$: $\dot{\alpha}_A = \gamma C_A \geq 0$ (requires $C_A > 0$).
  \textbf{Needs explicit Assumption $C_A > 0$.}
\item $\alpha_A = 1$: $\dot{\alpha}_A = -\zeta_\alpha R_A^{\text{surface}} \leq 0$.
  Correct as stated.
\item $\hat{M}_A, \Xi_A$ bounds: Mean-reverting dynamics ensure $[0,1]$
  invariance. Correct as stated.
\end{itemize}

\paragraph{Global existence.}
Since $\mathcal{X}$ is compact when $\Delta_{\max} < \infty$, and the
vector field points inward (or trajectories remain on the invariant
manifold) at every boundary, solutions cannot escape to infinity.
The escape lemma then implies $T_{\max} = \infty$. $\square$
\end{proof}

% --- Boundary Cases ---
\begin{boundary-case}
\textbf{$\sigma_A(0) = 0$ exactly:} The entire trajectory lies on the
invariant manifold $\sigma_A(t) \equiv 0$. This is the $\sigma$-trap
characterized in Theorem~4.1. The agent never develops schema coherence.

\textbf{$\sigma_A(0) \in (0, 1)$ but small:} The exact integral
\eqref{eq:exact-integral} shows the trajectory grows multiplicatively
if the bracket term is positive and decays to zero if it is negative.
This is the bifurcation behavior formalized in Proposition~4.1.

\textbf{Condition \eqref{eq:upper-bound-condition} fails:} If
$R_0 \geq 1/(1 - \gamma_\sigma)$, the analytical equilibrium exceeds 1.
The ODE alone cannot keep $\sigma_A \leq 1$. The numerical clamp in
Algorithm~3.1 is the effective enforcement mechanism, making the
implemented system a projected dynamical system rather than a pure ODE.

\textbf{$C_A = 0$ for an interval:} The $\alpha_A = 0$ boundary is not
reflecting; $\alpha_A$ remains at 0. This is physically interpretable as
attention collapse.
\end{boundary-case}

% --- Circular Dependency Check ---
\begin{circularity-check}
\textbf{Does Proposition~3.2 depend on any later result?}
The proof of $\sigma_A$ non-negativity via \eqref{eq:exact-integral}
depends only on the ODE structure and the assumption $\sigma_A(0) \geq 0$.
It does not require Proposition~3.6 (estimator consistency) or
Proposition~4.1 (bifurcation analysis). The upper bound discussion
references the equilibrium value $\sigma_A^*$ from Proposition~3.4,
but this is for interpretation only --- the forward invariance claim
itself is independent of Proposition~3.4.

\textbf{Does any later result depend on Proposition~3.2?}
Proposition~3.3 (fast-slow decomposition) requires that the state
space is compact and forward-invariant to apply Fenichel's theorem.
Proposition~3.4 (equilibria) requires the forward-invariant domain to
define equilibria. These are forward dependencies. Safe.

\textbf{Self-circularity:} The bound $\sigma_A \leq 1$ is not strictly
proved by the ODE alone for all parameter regimes. The proof must
be supplemented by the equilibrium analysis or the numerical clamp.
This is a \textbf{partial circular dependency}: Proposition~3.2's
claim of $\sigma_A \leq 1$ relies on Proposition~3.4's equilibrium
bound or Algorithm~3.1's clamping. We flag this for the reader.
\end{circularity-check}
